% TOP_PROBLEM: {"id": 3049, "title": "Square achievement game on an n x n grid", "category_id": 2, "category": {"id": 2, "name": "combinatorics", "display_name": "Combinatorics", "description": "Counting problems, graph theory, discrete structures.", "slug": "combinatorics", "order_index": 2, "created_at": "2026-07-31T15:26:25.671Z"}, "status": "open", "problem_number": "OPG-37228", "set_id": 12}
% FIRST_POSTED: 2026-10-01
% ATTEMPT_STATUS: unresolved
% UnsolvedMath ID 3049; source catalogue code OPG-37228
% !TeX program = lualatex
% Research attempt by GPT-6 Astra Ultra; no claim of a new complete solution.
\documentclass[11pt,a4paper]{article}
\usepackage[margin=25mm]{geometry}
\usepackage{fontspec}
\setmainfont{lmroman10-regular.otf}[BoldFont=lmroman10-bold.otf,
 ItalicFont=lmroman10-italic.otf,BoldItalicFont=lmroman10-bolditalic.otf]
\usepackage{amsmath,amssymb,mathtools}
\usepackage{unicode-math}
\setmathfont{latinmodern-math.otf}
\AtBeginDocument{\let\setminus\smallsetminus}
\usepackage{xurl,hyperref}
\hypersetup{colorlinks=true,urlcolor=blue}
\setcounter{secnumdepth}{0}
\setlength{\parindent}{0pt}
\setlength{\parskip}{5pt}
\setlength{\emergencystretch}{3em}
\begin{document}
\section*{Square achievement game on an n x n grid}\label{3049}
{\small UnsolvedMath ID 3049 · OPG-37228 · Combinatorics · OpenGarden}
\subsection{Definitions and mathematical statement}
For an integer $n\ge1$, the board is the set of cells
$[n]\times[n]$, initially empty. Two players, called $O$ and $X$, alternate
turns, with $O$ moving first. On each turn the player claims one previously
unoccupied cell by placing their own symbol there; a claimed cell never
changes owner, and passing is not allowed.

A winning set is
\[
 \{(i,j),(i+t,j),(i,j+t),(i+t,j+t)\},
 \qquad t\ge1,\quad 1\le i,j\le n-t.
\]
These are four cell positions at the corners of an axis-parallel square.
The first player to own all four members of any winning set wins immediately.
If every cell is filled and neither player has completed a winning set,
the game is drawn. Squares of every positive integral side length are
allowed; general rectangles and tilted squares are not winning sets.

Determine the outcome under optimal play for each $n$: a forced win for
$O$, a forced win for $X$, or a draw. A strategy specifies a legal move
for every game history where the player is to move. A forced win must
succeed against every sequence of legal opposing moves before the opponent
wins. A drawing strategy guarantees avoidance of defeat.

The existence of a completely filled board without a winning square only
shows that some play can end in a draw; it does not prove that either
player can force that play against optimal opposition. The task therefore
concerns strategies and the order of moves, not merely binary colourings.
\subsection{Short English statement}
Who can force a win, or a draw, when players alternately claim grid cells and try to complete the corners of an axis-parallel square?
\subsection{Research attempt}
\paragraph{Scope and process.}
This is a GPT-6 Astra Ultra research attempt dated 1 October 2026, using
primary-source browsing and elementary mathematical checks. The canonical
definition above is retained. Results below are restricted results or
reductions, not a claim of a new solution to the entire catalogue question.
No formal proof assistant or external mathematical referee checked this text.


\paragraph{Literature update and local scope.}
Jenrich's primary preprint [R1] reports drawing defences for the
second player on the $3\times3$ and $4\times4$ boards and a
winning strategy for the first player on $5\times5$. Therefore
the historical catalogue wording must not be interpreted as a
current claim that these outcomes have never been determined.
Here we independently prove the $n\le3$ outcomes and the general
strategy-stealing obstruction to a second-player forced win. We
do not certify the larger computer-assisted strategy trees in [R1].

\paragraph{Second cannot have a forced win on any board size.}
This is a finite strong positional game with the same winning sets
for both players; owning an extra cell cannot invalidate one's
winning set. Suppose the second player had a winning strategy $S$.
The first player initially claims an arbitrary spare cell, then
simulates playing second according to $S$ against the opponent's
actual moves. The simulated second-player cells are always a subset
of the first player's actual cells; the opponent's cells are the
same in both games. If $S$ requests the spare cell already owned,
mark it in the simulation and claim another free cell as the new
spare. Otherwise play the requested cell normally.

The opponent cannot win first, since the same sequence would defeat
$S$ in the simulated game. A simulated winning set is already an
actual winning set, so the first player wins by that stage or earlier.
If no free spare can be chosen, the actual board is full and the
simulation has only the old spare left; a guaranteed win by $S$
then lies already in the actual first player's cells, again implying
an earlier win. This contradiction rules out a second-player forced
win. Finite backward induction ensures that the remaining possibilities
are a first-player win or a draw; thus first can always avoid defeat.

\paragraph{An explicit second-player defence for $n=3$.}
Pair the following cells, leaving $(1,1)$ unpaired:
\[
 \{(1,2),(1,3)\},\quad
 \{(2,1),(2,2)\},\quad
 \{(2,3),(3,2)\},\quad
 \{(3,1),(3,3)\}.
\]
After each opposing move into a pair with empty partner, claim
that partner. If its partner is already owned by the defender,
or the unpaired cell was played, claim any available cell.
The invariant is that a pair never becomes wholly owned by the
opponent: the first opposing occupation is immediately answered
unless the partner was already protected. Extra defensive claims
cannot harm this invariant.

There are exactly five winning squares. The upper-right unit square
contains the first pair; both left unit squares contain the second;
the lower-right unit square contains the third; the side-two square
contains the fourth. Hence the first player can never complete a
square against this strategy. Coupled with first player's general
non-losing strategy, this proves a draw for $n=3$. For $n=1$ no
winning square exists. For $n=2$ the only winning set is the whole
board and neither player can own its four cells, so these are draws too.

\paragraph{Independent exact computation.}
The script
\texttt{scripts/check\_attempt\_batch\_algebra\_2026\_10\_01.py}
checks the pairing against all five squares and runs exact minimax
on the $3\times3$ board. A state records disjoint bitmasks of the
player to move and the opponent. A move completing a square scores
$+1$ immediately; otherwise its score is the negative of the
swapped successor state's value. A full board without a winner
scores zero. Maximizing over all legal moves is exhaustive backward
induction. The initial value computed is zero, consistent with the
explicit strategy proof. No larger-board search was run in this attempt.

\paragraph{Verification and first remaining task.}
The winning sets checked have equal horizontal and vertical side
lengths; allowing all rectangles would change the game and invalidate
the five-set calculation. A draw-colouring alone would not prove
a defence; the pairing gives legal responses to every history.
The unpaired cell and arbitrary-response cases are covered by the
invariant, including when the board becomes full. The remaining
independent task is to verify the reported $n=4,5$ strategies and
their application to larger boards while respecting immediate
opponent wins; this is not supplied merely by embedding a subboard.

\paragraph{Final status of this attempt.}
\textbf{UNRESOLVED as an independent complete strategy audit.}
Outcomes for $n\le3$ and exclusion of second-player forced wins
are proved. Reported stronger results are attributed to [R1], not
claimed as new discoveries or silently ignored.

\subsection{Sources}
\begin{itemize}
\item[S1] ULAM.AI and the UnsolvedMath Contributors, supplied problems.json, record 3049 (OPG-37228), OpenGarden collection. Catalogue identity and source transcription; CC BY 4.0.
\url{https://huggingface.co/datasets/ulamai/UnsolvedMath}.
\item[S2] Open Problem Garden, Square achievement game on an n x n grid. Original problem and discussion; the mathematical formulation below is expanded from the supplied source record.
\url{http://www.openproblemgarden.org/op/a_game_on_an_n_x_n_grid}.

\item[R1] Thomas Jenrich, \emph{Guaranteed successful strategies for
a square achievement game on an $n$ by $n$ grid}, arXiv:1109.2341,
2011. Primary preprint and description of the SQRGAME2 certificates;
abstract checked 1 October 2026.
\url{https://arxiv.org/abs/1109.2341}.

\end{itemize}
\end{document}
